\section{The Practical Task}

Describe what they have to build, and link the reference material:

\url{https://example.com/docs}

State the scope plainly -- which features are required, which are out of
scope, and what may be kept in memory rather than persisted.

\subsection{Steps}

\begin{enumerate}
    \item Read the reference documentation.
    \item Explain the approach before writing any code.
    \item Implement it.
    \item Analyse what you ended up with.
\end{enumerate}

\subsection{A code window}

\begin{neocode}[neocodetitlefont=\pixelfont]{config.yaml}
listen-on: 80

paths:
  "/":
    type: redirect
    target: "http://localhost:8080/index.html"
  "/dir/":
    type: servedir
    target: "/home/user/movies/"
\end{neocode}

\texttt{neocode} takes the title bar text as its mandatory argument and any
\texttt{listings} key in the optional one, e.g.
\texttt{[listing options=\{language=Python\}]}. Inside the block,
\texttt{@\ldots@} escapes back into \LaTeX{} so macros like \texttt{\textbackslash hwnumber}
still expand.

\subsection{A checklist}

\begin{neochecklist}
    \citem A checked item, written with \texttt{\textbackslash citem}
    \item An unchecked item, written with \texttt{\textbackslash item}
    \item Another thing to tick off before submitting
\end{neochecklist}
